% Part: sets-functions-relations
% Chapter: arithmetization
% Section: checking-details
%
\documentclass[../../../include/open-logic-section]{subfiles}


\begin{document}

\olfileid{sfr}{arith}{check}
\olsection{క్రమిత వలయాలు, క్షేత్రాలు}

ఈ అధ్యాయం అంతటా కొన్ని నిర్వచనాలు ``అవి ప్రవర్తించాల్సిన విధంగా''
ప్రవర్తిస్తాయని చెప్పాం. ఈ సాంకేతిక అనుబంధంలో దాని అర్థాన్ని పూర్తిగా
వివరించి, ఆ నిర్వచనాలు ``సరిగ్గా'' ప్రవర్తిస్తాయని (ఎలా చూపాలో
సంక్షిప్తంగా) తెలియజేస్తాం.

\olref[int]{sec}లో $\Int$పై సంకలనం, గుణకారం నిర్వచించాం. ఆ నిర్వచనాల
ప్రకారం అవి $\Int$కు మనం ``కోరుకునే'' నిర్మాణాన్ని ఇస్తాయని చూపాలి.
ప్రత్యేకంగా, అది ఒక వినిమయ వలయ నిర్మాణం.

\begin{defn}
	\emph{వినిమయ వలయం} అంటే, నిర్దిష్ట మూలకాలు $0$, $1$, అలాగే $+$,
	$\times$ క్రియలు అమర్చబడిన, ఈ ఎనిమిది సూత్రాలను సంతృప్తిపరిచే
	$S$ అనే సమితి:
	\begin{align*}
		\emph{Associativity}&&a + (b+ c) & = (a + b) + c \\
		&& (a \times b) \times c & = a \times (b\times c)\\
		\emph{Commutativity}&&a + b &= b+ a  \\
		&&  a \times b&= b\times a\\
		\emph{Identities}&&a + 0 &= a \\
		&& a \times 1 &= a\\
		\emph{Additive Inverse}&&(\exists b\in S)0&=a + b\\
		\emph{Distributivity}&&a \times (b+ c ) &= (a \times b) + (a \times c)
	\end{align*}
	అంతర్నిహితంగా, ఇవన్నీ $S$కు పరిమితమైన సార్వత్రిక పరిమాణీకరణలతో
	బంధించబడ్డాయి. ఇక్కడి $0$, $1$ మూలకాలు అదే పేర్లుగల సహజ సంఖ్యలు
	కానవసరం లేదని గమనించండి.
\end{defn}

కాబట్టి పూర్ణ సంఖ్యలు వినిమయ వలయాన్ని ఏర్పరుస్తాయని తనిఖీ చేయడానికి,
ఈ ఎనిమిది షరతులను సంతృప్తిపరుస్తామని మాత్రమే తనిఖీ చేయాలి. ఏ షరతునూ
స్థాపించడం {కష్టం} కాదు; కానీ ఇది కొంత శ్రమతో కూడినది. ఉదాహరణకు,
సంకలనం విషయంలో \emph{సహచర్యాన్ని} ఇలా నిరూపించవచ్చు:

\begin{proof}
$i, j, k \in \Int$లను స్థిరపరచండి. అప్పుడు $i =
\equivrep{a_1, b_1}{}$, $j = \equivrep{a_2,b_2}{}$, మరియు
$k = \equivrep{a_3, b_3}{}$ అయ్యే $a_1, b_1, a_2, b_2, a_3,
b_3 \in \Nat$ ఉంటాయి. (చదవడానికి సులభంగా, ``$\equivrep{x, y}{\Intequiv}$''కు
బదులుగా ``$\equivrep{x, y}{}$'' అని రాస్తాం; ఈ విభాగమంతా అలాగే
చేస్తాం.) ఇప్పుడు:
\begin{align*}
	i + (j + k) &= \equivrep{a_1, b_1}{}+(\equivrep{a_2, b_2}{} + \equivrep{a_3, b_3}{}) \\
	&= \equivrep{a_1,  b_1}{} + \equivrep{a_2+a_3, b_2+b_3}{}\\
	&= \equivrep{a_1 + (a_2 + a_3), b_1 + (b_2 + b_3)}{}\\
	&= \equivrep{(a_1 + a_2) + a_3, (b_1 + b_2) + b_3}{}\\
	&= \equivrep{a_1 + a_2, b_1 + b_2}{} + \equivrep{a_3, b_3}{}\\
	&= (\equivrep{a_1, b_1}{} + \equivrep{a_2, b_2}{}) + \equivrep{a_3, b_3}{}\\
	&= (i+j) + k
\end{align*}
ఇక్కడ $\Nat$పై సంకలనం ప్రవర్తనను స్వేచ్ఛగా ఉపయోగించుకున్నాం.
\end{proof}

అదే విధంగా, \emph{సంకలన విలోమాన్ని} ఇలా నిరూపించవచ్చు:

\begin{proof}
$i \in \Int$ను స్థిరపరచండి; ఏవో $a,b \in \Nat$లకు
$i = \equivrep{a,b}{}$. $j = \equivrep{b,a}{} \in \Int$గా
తీసుకోండి. సహజ సంఖ్యల ప్రవర్తనను ఉపయోగిస్తే,
$(a+b) + 0 = 0 + (a+b)$; అందువల్ల నిర్వచనం ప్రకారం
$\tuple{a+b, b+a} \sim_\Int \tuple{0,0}$, కాబట్టి
$\equivrep{a+b, b+a}{} = \equivrep{0, 0}{} = 0_\Int$. ఇప్పుడు
$i + j = \equivrep{a,b}{}+\equivrep{b,a}{}=
\equivrep{a+b, b+a}{}= \equivrep{0, 0}{} = 0_\Int$.
\end{proof}

\emph{వితరణ ధర్మానికి} నిరూపణ ఇదిగో:

\begin{proof}
పైన మాదిరిగా $i = \equivrep{a_1, b_1}{}$,
$j = \equivrep{a_2,b_2}{}$, మరియు $k = \equivrep{a_3, b_3}{}$లను
స్థిరపరచండి. ఇప్పుడు:
\begin{align*}
	i \times (j + k)
	&= \equivrep{a_1, b_1}{} \times (\equivrep{a_2,b_2}{} + \equivrep{a_3, b_3}{})\\
	&= \equivrep{a_1, b_1}{} \times \equivrep{a_2 + a_3,b_2+b_3}{}\\
	&= \equivrep{a_1  (a_2 + a_3) + b_1  (b_2+b_3), a_1  (b_2 + b_3) + b_1 (a_2 + a_3)}{}\\
	&= \equivrep{a_1 a_2 + a_1a_3 + b_1 b_2+b_1b_3, a_1 b_2 + a_1b_3 + a_2b_1 + a_3b_1}{}\\
%		&= \equivrep{a_1 a_2 + b_1b_2 + a_1a_3 + b_1 b_2+b_1b_3, a_1 b_2 + a_1b_3 + a_2b_1 + a_3b_1}{}\\
	&= \equivrep{a_1a_2 + b_1b_2, a_1b_2 + a_2b_1}{} + \equivrep{a_1a_3 + b_1b_3, a_1b_3 + a_3b_1}{}\\
	&= (\equivrep{a_1, b_1}{} \times \equivrep{a_2,b_2}{}) + (\equivrep{a_1, b_1}{} \times  \equivrep{a_3, b_3}{})\\
	&= (i \times j) + (i \times k)
\end{align*}
\end{proof}

మిగిలిన అయిదు షరతులను నిరూపించడాన్ని సాధనగా వదిలేస్తాం. అది చేసిన
తరువాత $\Int$ ఒక వినిమయ వలయమని, అంటే నిర్వచించిన సంకలనం, గుణకారం
అవి ప్రవర్తించాల్సిన విధంగానే ప్రవర్తిస్తాయని చూపినట్లవుతుంది.

\begin{prob}
$\Int$ ఒక వినిమయ వలయమని నిరూపించండి.
\end{prob}

కానీ మన పని ఇంకా పూర్తికాలేదు. $\Int$పై సంకలనం, గుణకారంతో పాటు
$\leq$ అనే క్రమ సంబంధాన్ని నిర్వచించాం; ఇది కూడా తగిన విధంగానే
ప్రవర్తిస్తుందని తనిఖీ చేయాలి. మరింత వివరంగా, $\Int$ ఒక
\emph{క్రమిత} వలయమని చూపాలి.\footnote{
\olref[sfr][rel][ord]{def:linearorder} నుంచి గుర్తు చేసుకోండి: సంపూర్ణ
క్రమం అనేది స్వావర్తన, సంక్రామక, ప్రతిసౌష్ఠవ, సంయుక్త ధర్మాలు గల
సంబంధం. క్రమ సంబంధాల సందర్భంలో సంయుక్తతను కొన్నిసార్లు
\emph{త్రివికల్పం} అంటారు; ఎందుకంటే ఏ $a$, $b$లకైనా
$a \leq b \lor a = b \lor a \geq b$.}

\begin{defn}
\emph{క్రమిత వలయం} అంటే, $\leq$ అనే సంపూర్ణ క్రమ సంబంధం కూడా
అమర్చబడి, కింది షరతులను సంతృప్తిపరిచే వినిమయ వలయం:
\begin{align*}
	a \leq b &\lif a + c \leq b + c\\
	(a \leq b \land 0 \leq c) &\lif a \times c \leq b \times c
\end{align*}
\end{defn}

\begin{prob}
$\Int$ ఒక క్రమిత వలయమని నిరూపించండి.
\end{prob}

మునుపటిలాగే, నిర్మించిన $\Int$ ఒక క్రమిత వలయమని చూపడం శ్రమతో
కూడినదే అయినా పరిపాటిగా చేయగలిగిన పని. దానిని మీకే వదిలేస్తాం.

దీనితో పూర్ణ సంఖ్యల పని ముగిసింది. ఇప్పుడు కరణీయ సంఖ్యల గురించి
చాలా సారూప్యమైన విషయాలను చూపాలి. ప్రత్యేకంగా, $+$, $\times$,
$\leq$లకు మనం ఇచ్చిన నిర్వచనాల కింద కరణీయ సంఖ్యలు ఒక క్రమిత
\emph{క్షేత్రాన్ని} ఏర్పరుస్తాయని చూపాలి:
\begin{defn}\ollabel{orderedfield}
\emph{క్రమిత క్షేత్రం} అంటే, ఈ షరతును కూడా సంతృప్తిపరిచే క్రమిత వలయం:
\begin{align*}
	\emph{Multiplicative Inverse}& & (\forall a \in S \setminus \{0\})(\exists b \in S) a\times b& = 1
\end{align*}
\end{defn}

$\Int$ ఒక క్రమిత వలయమని చూపిన తరువాత, $\Rat$ ఒక క్రమిత క్షేత్రమని
చూపడం సులభమే అయినా శ్రమతో కూడినది.

\begin{prob}
$\Rat$ ఒక క్రమిత క్షేత్రమని నిరూపించండి.
\end{prob}

పూర్ణ, కరణీయ సంఖ్యలను చూసిన తరువాత వాస్తవ సంఖ్యలను మాత్రమే చూడాలి.
ప్రత్యేకంగా, $\Real$ ఒక \emph{సంపూర్ణ} క్రమిత క్షేత్రం, అంటే సంపూర్ణతా
ధర్మం గల క్రమిత క్షేత్రం, అని చూపాలి. $\Real$కు సంపూర్ణతా ధర్మం ఉందని
\olref[cuts]{realcompleteness} స్థాపించింది. అయితే $\Real$ ఒక క్రమిత
క్షేత్రమని తనిఖీ చేసే (శ్రమతో కూడిన) పనిని ఇంకా పూర్తి చేయాలి.

\emph{ఆ} శ్రమతో కూడిన సాధనను మొదలుపెట్టేముందు, ఇంకొన్ని ``తక్షణ''
విషయాలను తనిఖీ చేయాలి. ఉదాహరణకు, ఏ కోతలు $\alpha$, $\beta$లకైనా
నిర్వచించిన $\alpha + \beta$ నిజంగా ఒక \emph{కోత} అని హామీ కావాలి.
ఆ వాస్తవానికి నిరూపణ ఇదిగో:

\begin{proof}
$\alpha$, $\beta$ రెండూ కోతలు కనుక
$\alpha + \beta = \Setabs{p + q}{p \in \alpha \land q \in \beta}$
అనేది $\Rat$ యొక్క శూన్యేతర నిజ ఉపసమితి. ఇప్పుడు ఏవో
$p \in \alpha$, $q \in \beta$లకు $x < p + q$ అనుకుందాం. అప్పుడు
$x - p < q$, కాబట్టి $x - p \in \beta$; అలాగే
$x = p + (x - p) \in \alpha + \beta$. అందువల్ల
$\alpha + \beta$, $\Rat$ యొక్క ఆద్య ఖండం. చివరగా, ఏ
$p + q \in \alpha + \beta$కైనా, $\alpha$, $\beta$ రెండూ కోతలు
కనుక $p < p_1$, $q < q_1$ అయ్యే $p_1 \in \alpha$,
$q_1 \in \beta$ ఉంటాయి; కాబట్టి
$p + q < p_1 + q_1 \in \alpha + \beta$; అందువల్ల
$\alpha + \beta$కు గరిష్ఠం లేదు.
\end{proof}

ఇలాంటి ప్రయత్నాలతో $\alpha - \beta$, $\alpha \times \beta$,
$\alpha \div \beta$ కూడా కోతలని తనిఖీ చేయవచ్చు (చివరి సందర్భంలో
$\beta$ సున్నా-కోత అయ్యే సందర్భాన్ని మినహాయించాలి). దీనినీ మీకే
వదిలేస్తాం.

\begin{prob}
$\Real$ ఒక క్రమిత క్షేత్రమని నిరూపించండి.
\end{prob}

కానీ ఒక చిన్న మిగులు అంశాన్ని పరిష్కరించాలి. \olref[cuts]{sec}లో
$\sqrt{2} = \Setabs{p \in \Rat}{p < 0 \text{ or }p^2 < 2}$గా
తీసుకోవచ్చని సూచించాం. అయితే ఈ సమితి ఒక \emph{కోత} అని చూపాలి.
ఆ వాస్తవానికి నిరూపణ ఇదిగో:

\begin{proof}
ఇది కరణీయ సంఖ్యల శూన్యేతర నిజ ఆద్య ఖండం అని స్పష్టం; కాబట్టి దీనికి
గరిష్ఠం లేదని చూపితే సరిపోతుంది. ప్రత్యేకంగా, $p$ అనేది
$p^2 < 2$ను సంతృప్తిపరిచే ధన కరణీయ సంఖ్య, అలాగే
$q = \frac{2p+2}{p+2}$ అయినప్పుడు $p < q$, $q^2 < 2$ రెండూ
అవుతాయని చూపితే సరిపోతుంది. $p < q$ అని చూడటానికి ఇవి గమనించండి:
\begin{align*}
	p^2 &< 2\\
	p^2 + 2p &< 2 + 2p\\
	p(p + 2) &< 2 + 2p\\
	p &< \tfrac{2+2p}{p+2} = q
\end{align*}
$q^2 < 2$ అని చూడటానికి ఇవి గమనించండి:
\begin{align*}
	p^2 &< 2\\
	2p^2 + 4p + 2 &< p^2 + 4p+ 4\\
	4p^2 + 8p + 4 &< 2(p^2 + 4p + 4)\\
	(2p+2)^2 & <2(p+2)^2\\
	\tfrac{(2p+2)^2}{(p+2)^2} &< 2\\
	q^2 &< 2
\end{align*}
\end{proof}
\end{document}
